% Ерөнхий дүгнэлт

\defformat
\chapter*{Ерөнхий дүгнэлт}
\addcontentsline{toc}{chapter}{Ерөнхий дүгнэлт}
\markboth{Ерөнхий дүгнэлт}{Ерөнхий дүгнэлт}
\myformat

Энэхүү дипломын ажлаар Монголын загварын сэтгүүлийн агуулгыг веб орчинд төвлөрүүлэн харуулах, нийтлэл, зураг, коллекц, дизайнер, чиг хандлага болон архивын мэдээллийг нэгдсэн бүтэцтэйгээр хэрэглэгчдэд хүргэх платформ боловсруулах зорилго тавьж, дараах үр дүнд хүрсэн.

Нэгдүгээр бүлэгт дижитал сэтгүүл ба онлайн мэдээний онолын үндэс, загварын медиа ба визуал редакцийн бүтэц, веб платформ хөгжүүлэх арга зүй, дотоод болон гадаадын төстэй платформуудын судалгаа, хэрэглэгчийн туршлага ба responsive дизайны үндэс, өгөгдөл ба контент зохион байгуулалтын зарчим, платформ хөгжүүлэхэд ашигласан технологиудын судалгааг авч үзсэн. Энэ судалгаа нь платформын онолын суурийг бүрдүүлсэн.

Хоёрдугаар бүлэгт Монголын загварын контентын одоогийн цахим орчны нөхцөл байдлыг шинжилж, Anoce/MFA платформын хамрах хүрээ, хэрэглэгчийн төрөл, функциональ болон функциональ бус шаардлага, use case, class, sequence, activity диаграмм, өгөгдлийн шаардлага, бизнес процессын шинжилгээг боловсруулсан.

Гуравдугаар бүлэгт шаардлагад үндэслэн платформын ерөнхий архитектур, дотоод бүтэц, модулийн зохиомж, API зохиомж, зохиомжийн шатны класс диаграмм, өгөгдлийн сангийн зохиомж, UI/UX wireframe, прототип, удирдлагын хэсэг, аюулгүй байдлын зохиомж, технологийн сонголтыг боловсруулсан.

Дөрөвдүгээр бүлэгт хөгжүүлэлтийн орчны бэлтгэл, frontend болон backend/API хэрэгжилт, өгөгдлийн сангийн хэрэгжилт, authentication/authorization, контент удирдлага, интерфейсийн хэрэгжилт, туршилт болон үр дүнгийн үнэлгээг тайлбарлав. Туршилтын үр дүнгээр платформын public болон admin үндсэн урсгалууд MVP түвшинд амжилттай ажилласан нь баталгаажсан.

Дипломын ажлын үр дүнд Монголын загварын сэтгүүл-архивын веб платформ хэрэгжиж, уншигч, бүртгэлтэй хэрэглэгч, редактор, админы үндсэн workflow-г дэмжих суурь бүрдсэн. Цаашид хэрэглэгчийн хувийн тохиргоо, semantic search, олон хэлний дэмжлэг, CDN болон кэшлэлт, автомат RAG синхрончлол зэрэг чиглэлээр платформыг сайжруулах боломжтой.
